\chapter{HASIL PEMBENTUKAN \textit{TRANSACTIONAL DATABASE}}
\label{chap:transaksi}\label{sec:hasil-t3}

\section{Pembentukan Dua \textit{Database}}
\label{sec:hasil-bentuk}
Dari 7.870.902 pasangan makalah-penulis terbentuk 1.821.757 transaksi pada \textit{author-based} dan 4.085.594 transaksi pada \textit{author-year-based}, dengan 4.373 \textit{venue} sebagai item unik pada keduanya. Gambar~\ref{fig:ilustrasi} mengilustrasikan perbedaannya pada satu penulis dengan 10 makalah di enam \textit{venue} (2021--2024): penulis ini menjadi satu transaksi berisi enam item pada \textit{author-based}, dan empat transaksi berisi tiga, tiga, satu, dan satu item pada \textit{author-year-based}. Dua makalah pada \textit{venue} yang sama dalam satu tahun, yaitu CCE pada 2023 dan ACC pada 2024, hanya menjadi satu item. Cuplikan berkas hasil ada pada Lampiran~\ref{lamp:transaksi}.

\begin{figure}[ht]
\centering
\figtikz{fig-ilustrasi-author}
\caption{Satu penulis pada \textit{author-based} dan \textit{author-year-based}}
\label{fig:ilustrasi}
\end{figure}

\section{Contoh Transaksi}
\label{sec:hasil-contoh}
Tabel~\ref{tab:contoh-a} dan Tabel~\ref{tab:contoh-ay} menunjukkan masing-masing sepuluh transaksi yang dipilih acak (\textit{seed} 42) menurut ukurannya. Transaksi berukuran empat sampai delapan item umumnya berisi \textit{venue} dari bidang yang berdekatan, sedangkan transaksi 25 item (\textit{Haibo Zhang}) mencampur banyak bidang dan kemungkinan menggabungkan beberapa orang bernama sama.

\begin{table}[ht]
\caption{Sepuluh contoh transaksi \textit{author-based}}
\label{tab:contoh-a}
\footnotesize
\begin{tabularx}{\textwidth}{@{}rlrX@{}}
\toprule
\textbf{tid} & \textbf{Author} & \textbf{Item} & \textbf{Isi transaksi} \\
\midrule
338077 & Cheoljoong Park & 1 & ISSCC \\
1245823 & Renann F. Brandao & 2 & LATINCOM, \textit{Image Perception, Observe}\ldots \\
851248 & Satoru Isaka & 3 & CASE, ICDL, SMC \\
466299 & Josef Guggemos & 4 & CSEDU, DeLFI, LTEC, RiE \\
389961 & Solvita Berzisa & 5 & BIR, ICEIS, ICIST, PROFES, SAC \\
808044 & Arun Kumar Thittai & 6 & HEALTHINF, EMBC, ISBI, \textit{Biomedical Applications i}\ldots, SMC, TENCON \\
64961 & Wei Li 0022 & 8 & ACAI, ECCV, ICDE, ICIP, ICPADS, NSDI, PCM, SIGIR \\
391158 & Thierry J. Chaussalet & 10 & BioInfoMed, CBMS, CIBCB, ECMS, HPCC/SmartCity/DSS, ICCS, ICIMTH, INISTA, \textit{IEEE Conf. on}\ldots \\
78397 & Tomás Masarík & 15 & ESA, GD, ICALP, IJCAI, ISAAC, IWOCA, IPEC [iwpec], MFCS, PODC, EC, \ldots{} ($+$5 lagi) \\
25570 & Haibo Zhang & 25 & WISA, ADHIP, AIAM, CACRE, CRC, CSPS, DSA, EEET, EIDWT, EMO, \ldots{} ($+$15 lagi) \\
\bottomrule
\end{tabularx}
\end{table}

\begin{table}[ht]
\caption{Sepuluh contoh transaksi \textit{author-year-based}}
\label{tab:contoh-ay}
\footnotesize
\begin{tabularx}{\textwidth}{@{}rlrrX@{}}
\toprule
\textbf{tid} & \textbf{Author} & \textbf{Tahun} & \textbf{Item} & \textbf{Isi transaksi} \\
\midrule
538906 & Hao Chen 0031 & 2018 & 1 & IJCNN \\
2566923 & Dirk Väth & 2023 & 2 & ACL, EACL \\
1626199 & Ayse Keles & 2024 & 3 & COMPSAC, EUVIP, EXPLIMED@ECAI \\
758263 & Rik Eshuis & 2015 & 4 & ICSOC, PRO-VE, NetWordS, OTM \\
627726 & Tom Mens & 2018 & 5 & BENEVOL, CDVE, ICSME, ICSR [icsr], MSR \\
1534896 & Christian Haubelt & 2016 & 6 & ARCS, CPS Data, EICS, MCSoC, SAMOS, WiMob \\
102172 & Lijie Wen 0001 & 2022 & 7 & COLING, EMNLP, ICASSP, IPDPS, KDD, NAACL-HLT, SIGIR \\
711013 & Syed Ishtiaque Ahmed & 2020 & 8 & CHI, COMPSAC, CSEE\&T, COMPASS, \textit{Graphics Interface}, ICTD, ICWSM, PDC \\
136050 & Javier Del Ser & 2020 & 10 & CEC, COMPSAC, CW, GECCO Companion, ICCS, IDEAL, IJCNN, ITSC, SOCO, SSCI \\
44703 & Feng Li & 2016 & 14 & CDC, CLSW, GECCO Companion, HCC, ICASSP, ICIC, IGARSS, INFOCOM, ISNN, IJCRS, \ldots{} ($+$4 lagi) \\
\bottomrule
\end{tabularx}
\end{table}

\section{Statistik dan Distribusi}
\label{sec:hasil-statistik-tx}
Tabel~\ref{tab:statistik-tx} merangkum statistik kedua \textit{database}, dan Gambar~\ref{fig:komposisi} menunjukkan distribusi item per transaksi.
\begin{enumerate}
\item \textit{Author-year-based} memiliki 2,24 kali lebih banyak transaksi, karena penulis yang aktif beberapa tahun menjadi beberapa transaksi, tetapi transaksinya lebih kecil (rata-rata 1,59 item dibanding 2,66).
\item Mayoritas transaksi hanya berisi satu item (58,4\% dan 71,0\%). Transaksi seperti ini tidak dapat membentuk pasangan, tetapi tetap masuk penyebut \textit{support}.
\item Distribusinya berekor panjang: rata-rata jauh di atas median (satu item), dan transaksi terbesar (332 dan 85 item) berasal dari \textit{Yang Liu}, nama umum yang menggabungkan banyak orang.
\end{enumerate}

\begin{table}[ht]
\caption{Statistik kedua \textit{transactional database}}
\label{tab:statistik-tx}
\small
\begin{tabular*}{\textwidth}{@{\extracolsep{\fill}}@{}lrr@{}}
\toprule
\textbf{Statistik} & \textbf{\textit{Author-based}} & \textbf{\textit{Author-year-based}} \\
\midrule
Jumlah transaksi & 1.821.757 & 4.085.594 \\
Item unik & 4.373 & 4.373 \\
Item per transaksi: median dan rata-rata & 1 dan 2,66 & 1 dan 1,59 \\
Item per transaksi: persentil 99 dan maksimum & 20 dan 332 & 7 dan 85 \\
Transaksi berisi 1 item & 58,4\% & 71,0\% \\
Transaksi berisi minimal 2 item & 41,6\% & 29,0\% \\
\bottomrule
\end{tabular*}
\end{table}

\begin{figure}[ht]
\centering
\figtikz{fig-komposisi-transaksi}
\caption{Komposisi ukuran transaksi pada kedua \textit{database}}
\label{fig:komposisi}
\end{figure}

\section{Perbedaan Kedua Representasi}
\label{sec:hasil-banding-repr}
Pada \textit{author-based}, dua \textit{venue} dianggap muncul bersama jika seorang penulis pernah menerbitkan di keduanya kapan saja pada 2015--2025. Pada \textit{author-year-based}, keduanya harus diterbitkan pada tahun yang sama. Tabel~\ref{tab:banding-repr} membandingkan pasangan antar-\textit{venue} dengan \textit{support} minimal 0,1\% pada \textit{author-based} (574 \textit{venue}).
\begin{enumerate}
\item \textit{Support} pasangan turun median 0,29 kali pada \textit{author-year-based}. Penyebabnya hampir seluruhnya penyebut: jumlah transaksi yang memuat pasangan hanya turun ke 0,65 kali, sedangkan jumlah transaksi \textit{database} naik 2,24 kali (0,65 dibagi 2,24 sama dengan sekitar 0,29). \textit{Min\_sup} yang sama karena itu berarti ambang absolut yang lebih tinggi.
\item Sekitar 23\% pasangan tidak pernah muncul pada \textit{author-year-based}, yaitu pasangan lintas tahun.
\item Urutan teratas relatif terjaga (Tabel~\ref{tab:top-pasangan}): 18 dari 20 pasangan teratas sama dan korelasi peringkatnya 0,892. ACL + EMNLP justru naik dari peringkat 6 ke 3 karena penulis yang konsisten menerbitkan di keduanya setiap tahun dihitung berulang.
\item Konferensi dua tahunan berselang tidak dapat berpasangan. ICCV dan ECCV memiliki 10.051 transaksi bersama pada \textit{author-based}, tetapi nol pada \textit{author-year-based} (Tabel~\ref{tab:iccv}).
\end{enumerate}

\begin{table}[ht]
\caption{Perbandingan pasangan \textit{venue} pada kedua representasi}
\label{tab:banding-repr}
\small
\begin{tabular*}{\textwidth}{@{\extracolsep{\fill}}@{}lr@{}}
\toprule
\textbf{Ukuran} & \textbf{Nilai} \\
\midrule
Pasangan yang muncul sedikitnya sekali, \textit{author-based} dan \textit{author-year-based} & 157.402 dan 120.937 \\
Median rasio \textit{support} (\textit{author-year} per \textit{author}) pada 100 pasangan teratas & 0,29 \\
Irisan 20 pasangan teratas dan korelasi peringkat Spearman & 18 dari 20 dan 0,892 \\
\bottomrule
\end{tabular*}
\end{table}

\begin{table}[ht]
\caption{Lima pasangan \textit{venue} teratas pada \textit{author-based} dan padanannya pada \textit{author-year-based}}
\label{tab:top-pasangan}
\footnotesize
\begin{tabularx}{\textwidth}{@{}>{\raggedright\arraybackslash}X rrrrrr@{}}
\toprule
\textbf{Pasangan} & \textbf{Trans. (A)} & \textbf{\textit{Support} (A)} & \textbf{Trans. (AY)} & \textbf{\textit{Support} (AY)} & \textbf{Rank AY} & \textbf{Rasio trans.} \\
\midrule
ICML + NeurIPS & 18.503 & 1,02\% & 17.719 & 0,43\% & 1 & 0,96 \\
ICLR + NeurIPS & 17.852 & 0,98\% & 16.438 & 0,40\% & 2 & 0,92 \\
AAAI + NeurIPS & 15.973 & 0,88\% & 12.281 & 0,30\% & 6 & 0,77 \\
CVPR + NeurIPS & 15.640 & 0,86\% & 14.485 & 0,35\% & 4 & 0,93 \\
CVPR + ICCV & 15.441 & 0,85\% & 12.359 & 0,30\% & 5 & 0,80 \\
\bottomrule
\end{tabularx}
\end{table}

\begin{table}[ht]
\caption{Transaksi \textit{author-year-based} yang memuat ICCV dan ECCV per tahun}
\label{tab:iccv}
\footnotesize
\setlength{\tabcolsep}{3.5pt}
\begin{tabular*}{\textwidth}{@{\extracolsep{\fill}}@{}lrrrrrrrrrrr@{}}
\toprule
 & \textbf{2015} & \textbf{2016} & \textbf{2017} & \textbf{2018} & \textbf{2019} & \textbf{2020} & \textbf{2021} & \textbf{2022} & \textbf{2023} & \textbf{2024} & \textbf{2025} \\
\midrule
ICCV & 1.543 & 0 & 1.959 & 0 & 3.627 & 0 & 5.971 & 0 & 8.439 & 0 & 15.260 \\
ECCV & 0 & 1.986 & 0 & 3.770 & 0 & 6.001 & 0 & 7.869 & 0 & 11.930 & 0 \\
\bottomrule
\end{tabular*}
\end{table}

Kedua representasi menjawab pertanyaan yang berbeda, dengan kelebihan dan keterbatasan masing-masing untuk \textit{frequent pattern mining}.
\begin{enumerate}
\item \textit{Author-based} berisi transaksi yang lebih kaya (41,6\% berisi minimal dua item) dan menangkap asosiasi lintas waktu, tetapi tidak memiliki dimensi waktu dan bebannya terkonsentrasi pada transaksi besar (0,1\% transaksi menanggung 27,8\% pasangan yang dapat dibentuk).
\item \textit{Author-year-based} mengikat kemunculan bersama pada waktu yang sama dan dapat dipotong per periode, tetapi 71,0\% transaksinya hanya satu item sehingga \textit{support} pasangan rendah dan pasangan lintas tahun hilang.
\item \textit{Author-based} dipilih untuk \textit{Task} 4 dan \textit{Task} 5 karena riwayat publikasi mencakup seluruh periode, dan \textit{author-year-based} dipilih untuk Bagian E karena memiliki dimensi tahun.
\end{enumerate}
Model himpunan tidak menyimpan jumlah makalah per \textit{venue}: 38,3\% makalah pada \textit{author-based} dan 17,4\% pada \textit{author-year-based} tidak menambah item baru.
